\chapter{Related Work}\label{relatedwork}
% In this chapter you demonstrate that you are sufficiently aware of the
% state-of-art knowledge of the problem domain that you have investigated as
% well as demonstrating that you have found a \emph{new} solution / approach / method.

Several approaches have been proposed to unify call-by-name and call-by-value evaluation strategies.
This section goes over the most notable ones.

Levy's approach \cite{levy2001thesis} introduces a programming paradigm called call-by-push-value (CBPV),
which is based on the quote ``a value is, a computation does''.
CBPV distinguishes between values and computations.
It can be constricted to simulate call-by-name in one way and call-by-value in another.
An extensive formalisation of CBPV in Rocq is provided by Forster et al. in \cite{ForsterSSS19}. 

Another approach to unifying \textsf{cbn} and \textsf{cbv} is by the use of linear logic. 
Maraist et al. \cite{MARAIST1999175} describe how two translations by Girard from intuitionistic logic into linear logic can be used to simulate call-by-name, call-by-value, and call-by-need evaluation strategies.
The idea is that linear logic can be used to keep track of the usage and evaluation of resources, which in this case are the arguments of a function.
Ehrhard and Guerrieri \cite{EhrhardGuerrieri2016} also used linear logic to unify \textsf{cbn} and \textsf{cbv}.
They introduced the Bang Calculus, which is a linear logic-based \lc that can express both \textsf{cbn} and \textsf{cbv} by using Girard's translations.
A formalisation of a linear logic-based λ-calculus in Rocq is provided by Zhao, Zhang, and Zdancewic in \cite{ZhaoZhangZdancewic2010},
with the implementation on GitHub \cite{ZdancewicLinearityRepo}.

The next approach we discuss uses thunks to unify \textsf{cbn} and \textsf{cbv}.
Thunks can be seen as wrappers around an expression that delays the evaluation until it is needed.
Continuous-Passing Style (CPS) translations are also a technique to unify \textsf{cbn} and \textsf{cbv}.
CPS translations are a way to rewrite a program such that all functions take an extra argument, which is the continuation function.
This continuation function is called with the result of the function, and it indicates the next step in the computation.
Hatcliff and Danvy \cite{HATCLIFF_DANVY_1997} used thunks and CPS translations to unify \textsf{cbn} and \textsf{cbv}.
There is one formalisation of a CPS translation in Rocq by Torrens \cite{torrens2024cpsrepo}, though it is not a direct translation of the approach by Hatcliff and Danvy.

The final approach uses model logic to unify \textsf{cbn} and \textsf{cbv}, and this is the approach we formalised in Agda.
Espírito Santo et al. \cite{espritosanto_et_al:LIPIcs:2019:10525} introduced the modal \lc \labox, which is a fragment of Bierman and de Paiva's intuitionistic modal logic (\textsf{IS4}) \cite{BiermanPaiva2000}.
The modal calculus \labox can express both \textsf{cbn} and \textsf{cbv} by using Girard's and Gödel's embeddings, respectively.
Espírito Santo et al. refined \labox to a simpler \lc, which they called the call-by-box \lc (\lab).
This refinement has simpler typing rules and a more straightforward reduction strategy.
No formalisation of \lab has been made in a proof assistant yet, and thus we formalised \lab in Agda.